\documentclass[11pt]{article}
\usepackage[T1]{fontenc}
\usepackage{lmodern}
\usepackage[a4paper,margin=25mm]{geometry}
\usepackage{amsmath,amssymb,amsthm}
\usepackage{microtype}
\usepackage[hidelinks]{hyperref}
\hypersetup{pdftitle={Positive initial coefficients and a positivity region for Brugidou's hook polynomials},
  pdfauthor={Xiaojian Zeng},pdfsubject={Partial coefficient positivity for hook partitions}}
\newtheorem{theorem}{Theorem}
\newtheorem{lemma}[theorem]{Lemma}
\newcommand{\Q}{\mathbb Q}
\newcommand{\Z}{\mathbb Z}
\newcommand{\N}{\mathbb Z_{\geq0}}
\newcommand{\coeff}[2]{[q^{#1}]#2}
\title{Positive initial coefficients and a positivity region\newline
for Brugidou's hook polynomials}
\author{Xiaojian Zeng\\[3pt]
\small PBC School of Finance, Tsinghua University, 100190\\
\small\texttt{xiaojianzeng@foxmail.com}}
\date{}

\begin{document}
\maketitle
\begin{abstract}
Brugidou associates an integer polynomial $J_\lambda(q)$ to each monomial
symmetric function under the specialization $e_j\mapsto q^{\binom j2}/j!$.
For hooks, write $K_{a,b}=J_{(a,1^b)}/b!$, where $a\geq2$ and $b\geq0$,
and put $d=\binom{a-1}{2}+(a-1)b$. We prove that $K_{a,b}\in\Z[q]$ and
that its coefficients in degrees $0$ through $\min(d,2b+3)$ are strictly
positive. All its coefficients are strictly positive if
$b\geq\max\{0,\lceil(\binom{a-2}{2}-4)/2\rceil\}$.
Together with the known row cases, this includes every hook with
$2\leq a\leq6$. The proof combines two coefficient-prefix identities with
a positive excess in the established row recurrence. Brugidou's
all-partition positivity conjecture remains unresolved by this result.
\end{abstract}

\section{Statement and context}

Let $\Lambda_{\Q}=\Q[e_1,e_2,\ldots]$ be the ring of symmetric functions,
where $e_j$ is the elementary symmetric function and $e_0=1$. For an
indeterminate $q$, define the unital specialization
\[
 \phi_q:\Lambda_{\Q}\longrightarrow\Q[q],\qquad
 \phi_q(e_j)=\frac{q^{\binom j2}}{j!}.
\]
For a nonempty partition $\lambda=(\lambda_1\geq\cdots\geq\lambda_r)$ of
$n$, put $\nu(\lambda)=\sum_{i=1}^r(i-1)\lambda_i$ and
$a_j(\lambda)=\#\{i:\lambda_i=j\}$. The monomial symmetric function
$m_\lambda$ is the sum of the distinct monomials with nonzero exponent
multiset $\lambda$. Brugidou's normalization is
\begin{equation}\label{eq:normalization}
 \phi_q(m_\lambda)=
 \frac{(1-q)^{n-r}q^{\nu(\lambda)}}{(n-1)!\prod_{j\geq1}a_j(\lambda)!}
 J_\lambda(q).
\end{equation}
The polynomial $J_\lambda$ belongs to $\Z[q]$, has nonzero constant
coefficient, and has degree
$\binom{n-1}{2}+r-1-\nu(\lambda)$
\cite[Corollary~3.5]{BrugidouMonomial}. All cancellations in
\eqref{eq:normalization} are identities in $\Q(q)$, not pointwise division
at $q=0$ or $q=1$.

Brugidou conjectures strict positivity at every degree from zero to
$\deg J_\lambda$ for every nonempty partition
\cite[Conjecture~4.1]{BrugidouMonomial}. We obtain the following restricted
result. Throughout, $J_m$ abbreviates $J_{(m)}$.

\begin{theorem}\label{thm:main}
For integers $a\geq2$ and $b\geq0$, define
\begin{align*}
 K_{a,b}&=\frac{J_{(a,1^b)}}{b!},&
 d_{a,b}&=\binom{a-1}{2}+(a-1)b,\\
 D_a&=\binom{a-2}{2},&
 F(a)&=\max\left\{0,\left\lceil\frac{D_a-4}{2}\right\rceil\right\}.
\end{align*}
Then $K_{a,b}\in\Z[q]$, and
\begin{equation}\label{eq:prefix-statement}
 \coeff{k}{K_{a,b}}>0
 \quad\text{for }0\leq k\leq\min(d_{a,b},2b+3).
\end{equation}
If $b\geq F(a)$, then $\coeff{k}{K_{a,b}}>0$ for every
$0\leq k\leq d_{a,b}$. Together with the established positive row
polynomials, this proves positivity for every hook with $2\leq a\leq6$.
\end{theorem}

The positive aggregate polynomials $J_n^{(r)}$ and their recurrence were
established in \cite[Theorem~6.3]{BrugidouAnalog}. Their two-variable
Tutte extension appears in \cite[Theorem~1.1]{BrugidouTutte}. These
aggregate statements concern sums over partitions of fixed size and
length, so positivity does not directly isolate an individual hook.
For individual non-column partitions, Brugidou proves a positive terminal block of
$n-1$ coefficients and reports computations through $n=43$
\cite[Corollary~5.2 and Section~4]{BrugidouMonomial}.
Our contribution is the initial range \eqref{eq:prefix-statement} and its
combination with a tail induction to obtain the region $b\geq F(a)$.
The elementary hook recurrence and the known row and column cases serve
as inputs. No equivalent hook statement was found in the primary sources
inspected for this note; this is a bounded comparison, not a priority claim.

Local coefficient data also suffice for Zeng's reconstruction of proper
trees with a multiplicity-isolated weighted core: the proof uses the
leading star-basis partition and its adjacent coefficient layer
\cite[Theorem~1.1 and Lemma~2.2]{ZengMultiplicity}. Our use of a coefficient
window concerns a different symmetric-function invariant and combines
initial positivity with a recursive tail identity.

\section{The hook recurrence and a row excess}

For a hook, \eqref{eq:normalization} reads
\begin{equation}\label{eq:hook-normalization}
 \phi_q(m_{(a,1^b)})=
 \frac{(1-q)^{a-1}q^{\binom{b+1}{2}}}{(a+b-1)!}K_{a,b}.
\end{equation}
It also gives $\deg K_{a,b}=d_{a,b}$. If $a\geq3$, the power sum
$p_{a-1}=\sum_i x_i^{a-1}$ satisfies
\[
 p_{a-1}e_{b+1}=m_{(a,1^b)}+m_{(a-1,1^{b+1})}.
\]
Indeed, its distinguished variable either occurs among the elementary
variables or does not; since $a-1\geq2$, each resulting monomial occurs
once. Specialization using \eqref{eq:hook-normalization} gives
\begin{equation}\label{eq:hook}
 (1-q)K_{a,b}=\binom{a+b-1}{a-2}J_{a-1}-q^{b+1}K_{a-1,b+1}.
\end{equation}
The base case from \cite[Equation~(3.22)]{BrugidouMonomial} is
\begin{equation}\label{eq:base2}
 K_{2,b}=1+q+\cdots+q^b.
\end{equation}
Since $K_{a,b}$ is already known to be a rational polynomial, induction on
$a$ in \eqref{eq:hook}, simultaneously for all $b\geq0$, proves its
integrality: division by $1-q$ takes integral coefficient prefixes.
This proves the divisibility by $b!$ required in the theorem, without
assuming it in advance. Constants in the same recurrence give
\begin{equation}\label{eq:constant}
 K_{a,b}(0)=\frac{(a+b-1)!}{(b+1)!}>0.
\end{equation}
For $a=2$, this agrees with \eqref{eq:base2}; for $a\geq3$ it uses
$J_{a-1}(0)=(a-2)!$. The next base case follows directly:
\begin{equation}\label{eq:base3}
 K_{3,b}=(b+2)\sum_{k=0}^{b}q^k
       +\sum_{j=1}^{b+1}(b+2-j)q^{b+j}.
\end{equation}
Thus all coefficients are positive for $a=2,3$.

To state the needed row estimate, recall the aggregate notation. For
$n\geq r\geq1$, let $p_n^{(r)}=\sum_{|\lambda|=n,\,\ell(\lambda)=r}m_\lambda$.
The polynomials defined by
\[
 \phi_q(p_n^{(r)})=
 \frac{(1-q)^{n-r}q^{\binom r2}}{r!(n-r)!}J_n^{(r)}(q)
\]
have positive integer coefficients and constant term $(n-r)!$, with
$J_n^{(1)}=J_n$ \cite[Theorem~6.3(i)]{BrugidouAnalog}.
At $r=1$, their established recurrence is
\begin{equation}\label{eq:row}
 J_m=\sum_{j=1}^{m-1}\binom{m-1}{j}q^{\binom j2}J_{m-1}^{(j)}
 \qquad(m\geq2)
\end{equation}
\cite[Equation~(6.5)]{BrugidouAnalog}.

\begin{lemma}\label{lem:excess}
For $m\geq3$, let $R_m=J_m-(m-1)J_{m-1}$. Then
\[
 R_m\in\N[q],\qquad R_m(0)=0,\qquad
 \coeff{1}{R_m}=\frac{(m-1)!}{2}.
\]
\end{lemma}
\begin{proof}
Subtract the $j=1$ summand in \eqref{eq:row}. Every remaining summand is
nonnegative and has positive valuation. Only $j=2$ contributes to degree
one, giving $\binom{m-1}{2}(m-3)!=(m-1)!/2$.
\end{proof}

\section{Positive prefixes and completion by tails}

\begin{proof}[Proof of Theorem~\ref{thm:main}]
Integrality was proved above, and \eqref{eq:base2}--\eqref{eq:base3} settle
$a=2,3$. Fix $a\geq4$, $b\geq0$, and abbreviate
\[
 s=b+1,\qquad N=a+b-1,\qquad K=K_{a,b}.
\]
Set
\begin{align*}
 A&=\binom N{a-2}J_{a-1},& B&=K_{a-1,b+1},\\
 P&=\binom N{a-3}J_{a-2},& C&=K_{a-2,b+2}.
\end{align*}
Two applications of \eqref{eq:hook} give
\begin{equation}\label{eq:double}
 (1-q)K=A-q^sB,\qquad (1-q)B=P-q^{s+1}C.
\end{equation}
Write $p_i=\coeff{i}{P}$ and extend every coefficient sequence by zero
outside its support. The binomial identity
\[
 \binom N{a-2}(a-2)=(s+1)\binom N{a-3}
\]
and Lemma~\ref{lem:excess} imply that
\begin{equation}\label{eq:delta}
 \Delta:=A-(s+1)P=\binom N{a-2}R_{a-1}\in\N[q],\qquad
 \Delta_1=\frac{s+1}{2}p_0>0,
\end{equation}
where $\Delta_i=\coeff{i}{\Delta}$.

Put $c_k=\coeff{k}{K}$. For $0\leq k<s$, the first equation in
\eqref{eq:double} gives $c_k=\sum_{i=0}^k\coeff{i}{A}>0$.
For $k=s+t$, with $0\leq t\leq s$, the second equation gives
\[
 \sum_{j=0}^t\coeff{j}{B}=\sum_{i=0}^t(t-i+1)p_i.
\]
The shifted $C$ term has not yet appeared. Substituting this expression
and $A=(s+1)P+\Delta$ into the prefix of the first equation yields
\begin{equation}\label{eq:prefix}
 c_{s+t}=\sum_{i=0}^{s+t}\Delta_i
       +\sum_{i=0}^t(s-t+i)p_i
       +(s+1)\sum_{i=t+1}^{s+t}p_i.
\end{equation}
Every summand is nonnegative. If $t<s$, the term $(s-t)p_0$ is positive.
If $t=s$, the first sum contains $\Delta_1>0$. Thus $c_k>0$ through
$k=2s$.

At $k=2s+1$, the double prefix receives its first contribution from $C$,
namely $C_0=K_{a-2,b+2}(0)>0$. The analogous identity is
\begin{align}
 c_{2s+1}
 &=\sum_{i=0}^{2s+1}\Delta_i+\sum_{i=0}^{s+1}(i-1)p_i
       +(s+1)\sum_{i=s+2}^{2s+1}p_i+C_0\notag\\
 &\geq\Delta_1-p_0+C_0
   =\frac{s-1}{2}p_0+C_0>0.\label{eq:endpoint}
\end{align}
Here $s\geq1$, so strictness also holds at $s=1$.
For $a\geq4$, $d_{a,b}\geq3s\geq2s+1$; hence this endpoint lies in
the polynomial's degree range. This proves
\eqref{eq:prefix-statement}. In particular, the prefix argument has used
no positivity assumption on $B$ or on the nonconstant coefficients of $C$.

For full positivity, induct on $a$, simultaneously for $b\geq F(a)$.
The function $F$ is nondecreasing, so $b+1\geq F(a-1)$ and the induction
hypothesis makes every coefficient of $B$ strictly positive through its
degree $L$. Let $D=\deg A=D_a$. Evaluating the polynomial identity
\eqref{eq:double} at $q=1$ gives $A(1)=B(1)$. For $k\geq D$,
\begin{equation}\label{eq:tail}
 c_k=A(1)-\sum_{j=0}^{k-s}\coeff{j}{B}
     =\sum_{j>k-s}\coeff{j}{B},
\end{equation}
where a prefix with negative upper limit is empty. The last sum is
strictly positive for every $k\leq s+L-1=d_{a,b}$.
Finally,
\[
 b\geq F(a)\quad\Longrightarrow\quad D_a-1\leq2b+3,
\]
so the positive prefix reaches every index below $D$ and the tail covers
all remaining indices. This includes the boundary $D_a=2b+4$.

The values $F(2),\ldots,F(6)$ are $0,0,0,0,1$.
The known row case $K_{6,0}=J_6$ supplies the remaining value $b=0$ at
$a=6$, proving the final assertion.
\end{proof}

\section{Examples, exact checks, and limits}

The overlap in the subtraction of \eqref{eq:hook} is already visible at
$(a,b)=(5,1)$. Here the coefficient of $q^2$ equals
$10(6+6+3)-20=130$, and
\begin{align*}
 K_{5,1}={}&60+120q+130q^2+110q^3+80q^4+55q^5\\
          &+35q^6+20q^7+10q^8+4q^9+q^{10}.
\end{align*}
Separating the supports in \eqref{eq:hook} would require
$b+1\geq D_a$. Identities \eqref{eq:prefix}--\eqref{eq:endpoint} control
the intervening overlap and give the smaller sufficient bound $F(a)$.
For example, $F(14)=31$. The hook $(14,1^{31})$ has size $45$ and degree
$481$; the theorem covers all degrees $0$ through $481$. The terminal
block supplied by \cite[Corollary~5.2]{BrugidouMonomial} consists of
degrees $438$ through $481$, and the example exceeds the reported
size-$43$ computational range.

Two implementations use exact integers and rational numbers. The first
computes rows from connected labeled graph counts using
$J_m=T_{K_m}(1,q)$ \cite[Equation~(1.2)]{BrugidouTutte}, then applies
\eqref{eq:hook}. It checks 66 direct specializations through size 12 and
177 selected hook cases with $2\leq a\leq16$, including 88 cases in the
full positivity region. A separately implemented check obtains power sums
from
\begin{equation}\label{eq:log}
 \sum_{m\geq1}\frac{(-1)^{m-1}}{m}\phi_q(p_m)t^m
 =\log\left(\sum_{j\geq0}\frac{q^{\binom j2}}{j!}t^j\right).
\end{equation}
This follows by taking the formal logarithm of
$\sum e_jt^j=\prod_i(1+x_it)$.
It then uses the augmented-monomial recurrence of
\cite[Theorem~1]{MercaAugmented}: for
$\widetilde m_\lambda=(\prod_j a_j(\lambda)!)m_\lambda$,
\[
 \widetilde m_{(\lambda_1,\ldots,\lambda_r)}
 =p_{\lambda_r}\widetilde m_{(\lambda_1,\ldots,\lambda_{r-1})}
 -\sum_{i=1}^{r-1}
 \widetilde m_{(\lambda_1,\ldots,\lambda_i+\lambda_r,\ldots,\lambda_{r-1})},
\]
with parts reordered when necessary. This computes all 120 hooks of
sizes 2 through 16 directly from the defining specialization.
Comparison with an unrolled form of \eqref{eq:hook}, divided from its
leading term, also reproduces the 177 recorded cases and adds 28 cases
with $17\leq a\leq20$ and
$b\in\{0,1,2,5,F(a)-1,F(a),F(a)+1\}$.
Across these 205 cases, it verifies 3,243 prefix identities and 37,827
tail coefficient identities. These computations validate formulas on
explicit finite domains; the proof above supplies the universal result.
The two scripts and their exact parameter lists are retained in the
accompanying computational dossier.

The function $F$ is a sufficient bound, with no optimality assertion.
For larger $a$, this note does not establish full positivity for general
hooks with $0<b<F(a)$, or for general non-hook partitions. It proves no
coefficient log-concavity statement. The original all-partition conjecture
therefore remains unresolved here.

\paragraph{Computational and AI assistance.}
OpenAI Codex assisted with proof development, exact-checking code,
literature inspection, and manuscript preparation. A separate model
context reconstructed the argument and produced the second implementation.
That pass was an additional automated check, not external peer review.
The computations use Python integer arithmetic and \texttt{fractions.Fraction};
no floating-point signs or proof-assistant certificate are used.


\small
\raggedright
\bibliographystyle{plain}
\bibliography{references}

\end{document}
